%%%%%%%%%%%%%%%%%%%%%%%%%%%%%%%%%%%%%%%%%%%%%%%%%%%%%%%%%%%%%%%%%%%%%%%%%%%%%%
\section{Conclusion and Future Work}    \label{sec:conclusion}
%%%%%%%%%%%%%%%%%%%%%%%%%%%%%%%%%%%%%%%%%%%%%%%%%%%%%%%%%%%%%%%%%%%%%%%%%%%%%%

% https://www.cs.utexas.edu/~EWD/transcriptions/EWD03xx/EWD303.html

<<<<<<< Updated upstream
This paper introduced Taskgrind\footnote{\url{https://gitlab.inria.fr/ropereir/valgrind}}: a Valgrind tool for parallel program memory accesses analysis.
It targets parallel programs that can be abstracted to a segment graph.
We implemented a determinacy race analysis, as it sometimes hint correctness error of asynchronous task-based programs; while also addressing some pitfalls bound to heavyweight dynamic binary instrumentation (DBI).
As opposed to compile-time instrumentation (such as LLVM Sanitizers), heavyweight DBI ensures the instrumentation of every instructions of the program (even close-source) - aiming at no false-negative report due to non-instrumented segments.
However, it required a lot of engineering to filter-out the many false-positive that can arise.
We have illustrated a few in this paper such as the inherent non-determinacy of parallel runtime systems, memory allocators recycling and thread/segment-local accesses.

In the future, we would like to extend the support of Taskgrind to parallel programming model interfaces, move-out of Taskgrind and parallelize determinacy-race analysis (which is an embarassingly parallel algorithm) and eventually add more analysis.
In addition, the current design for memory recycling have some drawbacks: simply removing memory free-operations is not a satisfactory solution; we also need to support programming models built-in memory allocators.
Finally, we would like to investigate on parallel-correctness of GPU programming models.

=======
\romain{just vomitting ideas here, needs to be re-ordered}
- E. Dijkstra~\cite{EWD303} also mentionned that "\emph{I shall argue that we must tailor our programs to the proof requirements}" - which was not the path taken by simulation program over the past decades (one could argue that was for the sake of productivity).
- Taskgrind is a first step towards parallel program with strong correctness guarantees.

h unreliable for ensuring program correctness


This paper introduced Taskgrind: a determinacy race detector relying on heavyweight dynamic binary instrumentation (DBI).
It targets application codes relying on a well-defined parallel programming model semantic, that needs to be bridged with its segment graph interfaces.
Taskgrind currently embeds an OMPT tool with support for most OpenMP parallel constructs (for loops, tasks, dependences, targets, detach, ...).


Many work-arounds for memory (free to NOOP, suppress errors on memory allocated by the LLVM allocator)
We could replace the LLVM allocator similarly to the system allocator, actually release the memory, and add a concept of "memory liveness"

TODO: Cilk/QThreads

As opposed to compile-time instrumentation (such as LLVM Sanitizers), heavyweight DBA ensures no determinacy-race may be missed due to non-instrumented code.
However, it requires a lot of engineering to filter-out the many false-positive that can arise.
We have illustrated a few in this paper such as memory recycling, thread-local accesses, and inherent non-determinacy of parallel runtime systems.

Future work
	- Engineering work
		- add support for other sync. primitives (mutex)
		- move analysis out of Valgrind (Algorithm on determinacy race is O(n²) embarassingly parallel)

	- Research
		- GPU offloading in presence - in particular with OpenMP, could compile target tasks as 'CPU' tasks
						
>>>>>>> Stashed changes
